\subsection{Integral transport of the original flag lattices}
\label{sec:flag-lattice-transport}

We now construct the representative required by the weighted
flag count.  The transformations are in the common unipotent
group of the two flags.  The lattice, its congruence condition
and its integral defects are retained throughout.

Use the normal form, spaces and notation of
Section~\ref{sec:unipotent-stabilizer-covolume}.  In addition,
assume that $\Lambda$ is unimodular, that both flags have the
prescribed leading ideals, and that their last vectors are
$z$ and $cz$, with $c\bar c=1$ and $h(z,z)=-T$ for a
positive integer $T$.  Write
\[
 A_a=\mathfrak d^{-(r-a)},\qquad
 H_0=X_j+X'_j+Ez,\qquad \Pi=H_0\mathbin\perp
                 (I+\langle f_a:a\geq j\rangle_E).
\]
Thus $V=\Pi\perp W$ and $\dim_E\Pi=2r+1$.
Here $\mathfrak d$ is the extension of an ideal of
$\mathcal O_F$, as in the original construction.
Equivalently, in this everywhere finite-unramified extension,
it is enough that $\mathfrak d$ be fixed by conjugation:
equal exponents at the two primes above a split prime, and
the single exponent at an inert prime, descend to $F$.

The ordered pairing matrix
$B=(h(e_a,e'_b))_{a,b<j}$ is a direct sum of signed scalar
antidiagonal blocks on consecutive intervals.  A block on
$[a,b)$ has entries
\[
       B_{uv}=t\epsilon_v\,\mathbf1_{u+v=a+b-1},
                    \qquad \epsilon_v\in\{1,-1\},\ t\ne0.
 \tag{FT1}
\]
The scalar antidiagonal blocks occupy the prescribed consecutive
intervals, with column $v$ multiplied by its original sign
$\epsilon_v$.
The source's reference module is exactly
\[
 P_0=\sum_{a<r}A_ae_a+
       \sum_{a<j}A_ae'_a+\mathcal O z+
       \sum_{j\leq a<r}A_a^{-1}f_a.
\]
These are the original ordered pairing and reference module;
the exact defining interfaces are recorded in
\cite{flag-lattice-transport-source}.
Let $\delta_l$ be its leading $l$-minor, set
$a_l=\sum_{a<l}(r-a)$, and put
\[
 \mathfrak a_l=(\delta_l)\mathfrak d^{-2a_l},\quad
 \Delta=N\mathfrak a_j,\quad
 D_b=\prod_{\substack{1\leq l\leq j\\\delta_l\ne0}}
                         N\mathfrak a_l,\quad N_T=T^f.
\]
The empty products are $1$.  In particular, when $j=0$,
$\Delta=D_b=1$.

\begin{lemma}[The ideal controlling the integral normal form]
\label{lem:flag-reference-dual-ideal}
Every nonzero $\mathfrak a_l$ above is integral.  The ideal
\[
 \mathfrak c=(T)\prod_{\substack{1\leq l\leq j\\\delta_l\ne0}}
                             \mathfrak a_l\overline{\mathfrak a_l}
 \tag{FT2}
\]
is nonzero, integral and fixed by conjugation, and
\[
 N\mathfrak c=(D_bN_T)^2,\qquad
 \mathfrak cP_0\subset P_0^\#,
 \qquad \mathfrak cP_0^\#\subset P_0.
 \tag{FT3}
\]
The dual in (FT3) is taken inside $\Pi$.
\end{lemma}

\begin{proof}
The successive quotients of $\Lambda\cap X$ are the
specified ideals $A_a$.  Since fractional ideals over a
Dedekind domain are projective, these quotients can be split
successively.  Thus, for $a<j$, there are vectors
$u_a=e_a+\sum_{b<a}q_{ba}e_b$ such that
\[
                  \Lambda\cap X_j=\bigoplus_{a<j}A_au_a.
 \tag{FT4}
\]
The same assertion holds for the primed flag.  This does
not assert that $A_ae_a$ itself lies in $\Lambda$.
At a finite prime choose a generator $\varpi$ for the
localization of $\mathfrak d$ in the real field.  The
vectors $\varpi^{-(r-a)}u_a$ are a local lattice basis.
Their leading Gram determinant against the primed basis is
$\delta_l\varpi^{-2a_l}$, because both triangular changes
have determinant $1$.  Its entries are integral.  This
proves the integrality of $\mathfrak a_l$ at every prime,
under the actual successive-quotient hypothesis.

In these weighted coordinates a block in (FT1) is again
scalar antidiagonal, since $u+v$ is constant on its support.
If the block has length $b-a$, its scalar $\lambda$ satisfies
\[
 (b-a)v(\lambda)=
       v(\delta_b\varpi^{-2a_b})-
       v(\delta_a\varpi^{-2a_a}).
\]
Both endpoint valuations are nonnegative and at most
$v(\mathfrak c)$.  Consequently
$|v(\lambda)|\leq v(\mathfrak c)$.
The same argument applies to the conjugate block.  Thus
multiplication by $\mathfrak c$ makes both the weighted
pairing matrix on $X_j+X'_j$ and its inverse integral.
The pairing on the weighted $I,I^*$ summands is perfect.
The final block is $-T$, whose inverse denominator is
absorbed by $\mathfrak c\subset(T)$.  This proves (FT3).
The norm identity follows directly from (FT2), including
the rational degree factor $N(T)=T^{2f}=N_T^2$.
\end{proof}

Let $R$ denote the intersection of the two original flag
unipotent groups, and let $R_s$ be its subgroup fixing $W$
pointwise.  An admissible mod-$g$ flag congruence is also
preserved by $R$: pulling a lattice lift back by an element
of $R$ preserves its leading coefficient, and applying that
element again preserves the difference in $g\Lambda$.
This is the original congruence condition, not a choice of
a new graph.

\begin{theorem}[A nearby representative in the original orbit]
\label{thm:flag-lattice-near-representative}
There is $u\in R$ such that $\Lambda'=u\Lambda$ satisfies
\[
 \Lambda'\cap(W+I)=(\Lambda'\cap W)+(\Lambda'\cap I),
 \qquad
 \mathfrak c^{12r+50}P_0\subset\Lambda'\cap\Pi
                  \subset\mathfrak c^{-(12r+50)}P_0.
 \tag{FT5}
\]
The transformation giving the two inclusions can be chosen
in $R_s$ once the first, splitting condition holds.  All
successive flag ideals and the original mod-$g$ condition
are preserved.
\end{theorem}

The first condition in (FT5) implies the source's
\texttt{Split} condition.  Its theta lattice is
$\bigl(\Lambda'\cap(W+I)\bigr)+I$; after this splitting
it is $(\Lambda'\cap W)+I$, as required by that definition.

\begin{proof}
We give the construction, with the smaller exponent $6r+10$.

\smallskip\noindent\emph{Splitting the definite quotient.}
The module $\Lambda\cap I$ is saturated in
$\Lambda\cap(W+I)$, and its quotient is projective.
Choose a module splitting and take the $E$-span of its
image.  This gives a complement $W'$ to $I$ in $W+I$,
with the same integral quotient lattice.  Write
$W'=\{w+a(w):w\in W\}$, where $a:W\to I$ is linear.
A unitary transformation fixing $V_0$ can send this graph
to $W$ and act identically on $(W+I)/I$.
For completeness, if $A:W\to I$ is linear and $A^*:I^*\to W$
is its adjoint for the perfect $I,I^*$ pairing, the matrix
on $I\oplus W\oplus I^*$
\[
 n_A=
 \begin{pmatrix}
  1&A&-w_0AA^*\\
  0&1&-A^*\\
  0&0&1
 \end{pmatrix},\qquad w_0+\bar w_0=1,
 \tag{FT6}
\]
is unitary.  It is the identity on $X_j+X'_j+Ez$.
Taking $A=-a$ proves the assertion.  The pairing identities
in (FT6) use only $w_0+\bar w_0=1$.  Thus this transformation
is in the original $R_1\subset R$, identified above.

\smallskip\noindent\emph{Simultaneous control of the paired flags.}
Choose the two splittings (FT4), with upper unitriangular
basis matrices $U,U'$.  Their integral Gram matrix is
\[
                         G=U^tB\overline{U'}.
\]
There is a single common flag isometry making every
diagonal block of $U^t$ the identity.  In fact let $S$ be
block diagonal, with upper unitriangular blocks equal to
the inverses of the corresponding blocks of $U$, and put
\[
                    S'=\overline{B^{-1}(S^t)^{-1}B}.
 \tag{FT7}
\]
Conjugation by a signed antidiagonal block reverses the
triangular order, so $S'$ is upper unitriangular on the
same blocks.  The identity $S^tB\overline{S'}=B$ proves
that these two actions extend to one unitary map, equal
to the identity on $I+I^*+Ez+W$.  It lies in $R_s$.
This verifies simultaneity even when $j=r$.

After applying (FT7), write the unchanged matrix as
$G=ABC$, where $A$ is lower unitriangular with identity
diagonal blocks, and $C$ is upper unitriangular.  Work
locally in the weighted coordinates used above, so $G$
is integral.  For a column in a block ending at $l$,
replace that row of the leading $l\times l$ minor of
$G$ by any later row.  The determinant is, up to sign,
$\det(B_l)$ times the corresponding entry of $A$.
Indeed the relevant column of $A_l^{-1}$ is the standard
basis column, since the last diagonal block of $A_l$ is
the identity.  Thus $\mathfrak cA$ is integral.

For a block beginning at $l'$, elimination of the preceding
rows gives the Schur complement
\[
 G_{\mathrm{block},*}
  -G_{\mathrm{block},<l'}G_{<l',<l'}^{-1}G_{<l',*}
                       =B_{\mathrm{block}}C_{\mathrm{block},*}.
 \tag{FT8}
\]
Its entries have denominators dividing the preceding
leading determinant, which is absorbed by $\mathfrak c$.
By (FT3), the inverse block costs at most one further
power of $\mathfrak c$.  Hence $\mathfrak c^2C$ is
integral.  This argument is also valid for the first
block, with empty Schur complement.

An upper or lower unitriangular $j\times j$ matrix has
inverse $\sum_{a=0}^{j-1}(-N)^a$.  Therefore the two
transformed flag lattices and their reference lattices
contain each other after multiplication by
$\mathfrak c^n$, where $n=2r+2$.  All statements are
local ideal inclusions, so they hold globally without a
choice of a global generator of $\mathfrak d$ or
$\mathfrak c$.

\smallskip\noindent\emph{The common isotropic part and the last vector.}
Let $U_X=\Lambda\cap X$ and $U'_j=\Lambda\cap X'_j$,
using the transformed lattice.  Projection along $X_j$
maps $U_X$ to a lattice $\widetilde I\subset I$ with
successive ideals $A_j,\ldots,A_{r-1}$.  The pairing
between $\Lambda\cap X_j$ and $U'_j$ has determinant
ideal $\mathfrak a_j$ or its conjugate.  Its adjugate
therefore shows
\[
                      \mathfrak c\widetilde I
                           \subset\Lambda\cap I
                           \subset\widetilde I.
 \tag{FT9}
\]
Explicitly, if $x+i\in U_X$ with $x\in X_j$, then
integrality against $U'_j$ places $x$ in its paired
dual; multiplication by $\mathfrak c$ puts $x$ in
$\Lambda\cap X_j$.  Subtraction proves (FT9).

Split the flag filtration of $\widetilde I$ projectively.
An upper unitriangular map $D$ on $I$ carries it to
$P_I=\sum_{a\geq j}A_ae_a$.  Extend $D$ to the dual
space by the inverse adjoint and fix $H_0+W$.  This is
again a single element of $R_s$ and does not alter the
paired flag bounds.  Now
\[
                  \mathfrak cP_I\subset\Lambda\cap I\subset P_I.
 \tag{FT10}
\]
The last leading coefficient gives a lift
$z+x+i\in\Lambda$, with $x\in X_j$ and $i\in I$.
The same adjugate argument gives
$\mathfrak c(z+i)\subset\Lambda$.  The Eichler
transformation (US11), with isotropic vector $i$ and
$u=z/T$, sends $z+i$ to $z$, fixes $I$, and fixes
$X_j+X'_j+W$.  It lies in $R_s$.  We have reached
\[
 \mathfrak c^n A_0\subset\Lambda,
 \qquad A_0=P_0\cap V_0=P_I+(P_0\cap H_0),
 \qquad \Lambda\cap I\subset P_I.
 \tag{FT11}
\]
Every transformation used since the first step fixes $W$,
so the splitting condition is preserved.

\smallskip\noindent\emph{The dual isotropic directions.}
Let $P_{I^*}=\sum_{a\geq j}A_a^{-1}f_a$.
The split module
$M+(\Lambda\cap I)=\Lambda\cap(W+I)$ is saturated
in $\Lambda$.  For $y\in P_{I^*}$ prescribe on it the
Hermitian functional equal to $h(y,\cdot)$ on
$\Lambda\cap I$ and zero on $M$.
It is integral by (FT10).  Projectivity and unimodularity
extend it to pairing with an element $\Phi y\in\Lambda$.
The extension can be chosen $\mathcal O$-linearly in $y$,
since $P_{I^*}$ is projective.  Extend it to an $E$-linear
map.  Its values are perpendicular to $W$, and
\[
       \Phi y=y+a(y)+b(y),\qquad
       a(y)\in H_0,\quad b(y)\in I,
       \qquad\Phi(P_{I^*})\subset\Lambda\cap\Pi.
 \tag{FT12}
\]

Write $P_H=P_0\cap H_0$.  Pairing (FT12) with
$\mathfrak c^nP_H\subset\Lambda$ gives
$a(P_{I^*})\subset\mathfrak c^{-(n+1)}P_H$ by (FT3).
Pairing two values of $\Phi$ then shows
\[
        (b+b^*)(P_{I^*})\subset\mathfrak c^{-(2n+3)}P_I.
 \tag{FT13}
\]
Here $b^*$ is the adjoint for the perfect $I,I^*$ pairing.
Indeed the Hermitian form of $b+b^*$ is
$h(\Phi y,\Phi y')-h(a(y),a(y'))$; the first term is
integral, and the second has denominator at most
$\mathfrak c^{2n+3}$, because
$h(P_H,P_H)\subset\mathfrak c^{-1}$.

The map $K=b-w_0(b+b^*)$ is skew adjoint.  The
transformation that sends $y\in I^*$ to $y-Ky$ and
fixes $V_0+W$ is unitary and belongs to $R_s$.
Its transformed lift is
$y+a(y)+w_0(b+b^*)y$.  Equations (FT11)--(FT13) imply
\[
       \mathfrak c^{3n+3}P_0\subset\Lambda'\cap\Pi.
\]
Conversely integrality of $\Lambda'$ and (FT3) give
$\Lambda'\cap\Pi\subset\mathfrak c^{-(3n+4)}P_0$.
Thus the two-sided exponent $3n+4=6r+10$ suffices, and
implies (FT5).  If $j=r$ the dual isotropic module is
zero and this last step is vacuous.  If $j=0$ the
paired-flag step is vacuous.  The construction also
applies when $r=1$ or $r=0$.
\end{proof}

\begin{theorem}[The two remaining stabilizer quotients]
\label{thm:flag-small-stabilizer-quotients}
Suppose the original lattice is split and
$\mathfrak c^eP_0\subset\Lambda\cap\Pi
\subset\mathfrak c^{-e}P_0$ for an integer $e\geq0$.
Let $\Gamma_B=R\cap\operatorname{Stab}(B)$, let $R_1$
fix $V_0$, and let $Z_I\subset R_1$ fix $W$.
Each of the source's central and top covolume ratios
is at most
\[
                     (N\mathfrak c)^{(3e+11)(2r+1)^2}.
 \tag{FT14}
\]
Together with (US7), this proves the original bound
\[
 \operatorname{ivol}(\Gamma_{\Lambda_{\mathrm{ref}}},\Gamma_\Lambda)
 \leq b^k\Delta^r
                  (N\mathfrak c)^{2(3e+11)(2r+1)^2},
 \tag{FT15}
\]
whenever the original index constraint $xi\leq\Delta$
holds.
\end{theorem}

\begin{proof}
Put $P=\Lambda\cap\Pi$.  The inclusions and (FT3) give
$P^\#\subset\mathfrak c^{-(e+1)}P_0$.
If a unitary map $s$ fixes $W$ and
\[
                      (s-1)P_0\subset\mathfrak c^{2e+1}P_0,
 \tag{FT16}
\]
then $(s-1)P^\#\subset P$.
The $\Pi$-projection of $\Lambda$ is contained in $P^\#$,
so $(s-1)\Lambda\subset P\subset\Lambda$.
Unitarity makes the additive index $[\Lambda:s\Lambda]$
equal to $1$.  Thus $s\Lambda=\Lambda$.

For the central quotient, elements of
$\Gamma_{\Lambda_{\mathrm{ref}}}\cap Z_I$ preserve
$P_0=\Lambda_{\mathrm{ref}}\cap\Pi$.
The kernel of reduction on
$P_0/\mathfrak c^{3e+11}P_0$ satisfies (FT16) and
lies in $\Gamma_\Lambda\cap Z_I$.
For a projective rank-$q$ module, the number of
endomorphisms modulo an ideal $\mathfrak a$ is
$(N\mathfrak a)^{q^2}$: localize and use a basis,
then combine the finite primary components.
Taking $q=2r+1$ proves the central bound.

For the top quotient, work on
$A_0=\Lambda_{\mathrm{ref}}\cap V_0=P_I\oplus P_H$.
An element $s\in R$ induces on $I\oplus H_0$ a
matrix
\[
             \begin{pmatrix}D&B\\0&A\end{pmatrix},
\]
with $A$ unitary and $D$ upper unitriangular on the
isotropic flag.  Suppose $(s-1)A_0\subset\mathfrak c^NA_0$,
where $N=3e+11$.  Then $A,D$ preserve their respective
lattices, as do their inverses: their additive determinant
norms are $1$.  Moreover their differences from the
identity lie in $\mathfrak c^N$.
Put $C=BA^{-1}:H_0\to I$.  It sends $P_H$ to
$\mathfrak c^NP_I$, while its adjoint sends $P_{I^*}$
to $\mathfrak c^{N-1}P_H$, by (FT3).

Let $n_C$ denote (FT6), formed with $H_0$ and $C$.
The ordered product
\[
                    n_{BA^{-1}}
                       \operatorname{diag}(D,A,D^{-*})
\]
gives a unitary lift of the action of $s$ on $V_0$,
fixing $W$: its upper right block on $I\oplus H_0$
is $CA=B$.
The quadratic entry $-w_0CC^*$ lies in
$\mathfrak c^{2N-1}$, so this lift differs from the
identity on $P_0$ by an element of
$\mathfrak c^{N-1}P_0$.  It belongs to $R_s$, because
it has the same action as $s$ on both flags and fixes
$W$.  Since $N-1\geq2e+1$, (FT16) makes it preserve
$\Lambda$.  Its action on $V_0$ is therefore in the
image of $\Gamma_\Lambda$.

Reduction on $A_0/\mathfrak c^NA_0$ consequently bounds
the top relative index by
$(N\mathfrak c)^{N(\dim V_0)^2}$, at most (FT14).
The inverse covolume ratio is bounded by that relative
index.  Finally $Z_I\subset R_1\subset R$ are normal,
because $R$ preserves $V_0$ and sends $W$ into $W+I$;
thus it preserves $V_0+W$, whose pointwise stabilizer
in $R$ is $Z_I$.  The index formula along this chain factors the full
covolume ratio into the central, middle and top ratios.
Apply (US7) and $(xi)^{r-j}\leq\Delta^r$ to obtain
(FT15).
\end{proof}
